% TOP_PROBLEM: {"id": 3100055, "title": "Suppose that S is a set of positive integers with the property that S+S -- that is, all sums of the form s_(1)+s_(2) for", "category_id": 1, "category": {"id": 1, "name": "number_theory", "display_name": "Number Theory", "description": "Properties of integers, prime numbers, Diophantine equations.", "slug": "number-theory", "order_index": 1, "created_at": "2026-07-31T15:26:25.670Z"}, "status": "partially_solved"}
% FIRST_POSTED: null
% ENTRY_KIND: statement_only
% Imported from: batch11.tex; UnsolvedMath ID 3100055
\documentclass[11pt]{article}
\usepackage{amsmath,amssymb,mathtools}
\usepackage{fontspec,unicode-math}
\setmainfont{Latin Modern Roman}
\setmathfont{Latin Modern Math}
\usepackage{xurl,hyperref}
\newcommand{\sref}[2]{\hyperlink{source:#1:#2}{[S#2]}}
\newcommand{\cataloguescope}[1]{\paragraph{Mathematical statement.}#1}
\begin{document}
% UnsolvedMath ID 3100055; source catalogue code AMR-030-0055
\setcounter{section}{3100054}
\section{Suppose that S is a set of positive integers with the property that S+S -- that is, all sums of the form s\_(1)+s\_(2) for}\label{3100055}
{\small\sffamily UnsolvedMath ID 3100055 · Number Theory}
\subsection{Definitions and mathematical statement}

\paragraph{Shared notation.}$\mathbb N=\{1,2,\ldots\}$, and $\mathbb Z,\mathbb Q,\mathbb R,\mathbb C$ denote the integers, rationals, reals and complex numbers. Any other conventions are specified locally.


For $S\subset\mathbb N$, put $S+S=\{a+b:a,b\in S\}$ and
\[r_S(n)=\#\{(a,b)\in S^2:a\le b,\ a+b=n\}.\]
Repeated summands are allowed. The set $S$ is an asymptotic additive basis of order two if $r_S(n)\ge1$ for every sufficiently large integer $n$. Counting ordered representations instead changes this count by at most a factor of two, so boundedness is unchanged.

\paragraph{Statement recorded in the supplied source.}
Must every asymptotic additive basis $S\subset\mathbb N$ of order two satisfy
\[\limsup_{n\to\infty}r_S(n)=+\infty?\]
Equivalently, can there be such an $S$ and a finite constant $C$ with $r_S(n)\le C$ for every $n\in\mathbb N$? The conjectured answer to the existence question is negative. \sref{3100055}{2}

\subsection{Short English statement}
Must the number of ways to write an integer as the sum of two elements of an asymptotic additive basis become arbitrarily large?
\subsection{Sources}
\begin{description}
\item[\hypertarget{source:3100055:1}{S1}] ULAM.AI and the UnsolvedMath Contributors, \emph{UnsolvedMath: A Curated Collection of Open Mathematics Problems}, record 3100055 (AMR-030-0055). CC BY 4.0. \url{https://huggingface.co/datasets/ulamai/UnsolvedMath}.
\item[\hypertarget{source:3100055:2}{S2}] Melvyn B. Nathanson, \emph{Generalized additive bases, König\textquotesingle s lemma, and the Erdős--Turán conjecture} (2003), Section1, unordered representation function, asymptotic bases, and displayed Erdős--Turán conjecture, pages1--2. \url{https://arxiv.org/abs/math/0302155}.
\end{description}
\label{end:3100055}
\end{document}
